\documentclass[11pt,a4paper]{article}

\usepackage[margin=2.5cm]{geometry}
\usepackage{amsmath}
\usepackage{booktabs}
\usepackage{xcolor}
\usepackage{listings}
\usepackage[colorlinks=true,linkcolor=blue,urlcolor=blue,citecolor=blue]{hyperref}
\usepackage{parskip}

\lstdefinestyle{deck}{
  basicstyle=\ttfamily\small,
  breaklines=true,
  frame=single,
  columns=fullflexible,
  keepspaces=true
}
\lstdefinestyle{python}{
  language=Python,
  basicstyle=\ttfamily\small,
  breaklines=true,
  frame=single,
  columns=fullflexible,
  keepspaces=true,
  commentstyle=\color{gray},
  showstringspaces=false
}

\title{Custom Laser Profile Injection}
\author{Zheyuan Chen \\ York Plasma Institute, University of York}
\date{\today}

\begin{document}
\maketitle

This document describes the modifications made to EPOCH's \texttt{laser}
block to support loading arbitrary laser amplitude and phase profiles from
raw binary data files\footnote{\url{https://epochpic.github.io/documentation/input_deck/binary_files.html}}.
See the EPOCH documentation\footnote{\url{https://epochpic.github.io/documentation/input_deck/input_deck_laser.html}}
for the standard \texttt{laser} block parameters that are not covered here.

The custom profile injection feature allows a laser field of the form
$E = A \cdot \sin(\omega t + \varphi)$ --- with $A$ and $\varphi$ arbitrary
functions of the transverse boundary coordinates and (optionally) time ---
to be driven entirely from externally generated data files, bypassing
EPOCH's built-in analytic profile functions. This is particularly useful
when interfacing with external beam propagation codes such as
LASY\footnote{\url{https://github.com/LASY-org/lasy}}, or when the target
profile cannot be expressed as one of EPOCH's analytic primitives. A Python
wrapper for LASY written by the same author can be found
here\footnote{\url{https://github.com/ZheyuanChen/AELP/tree/main}}.

The feature is implemented in \textbf{epoch2d} and \textbf{epoch3d}. epoch1d is unmodified.

Two modes exist, selected per laser block:

\begin{itemize}
  \item \textbf{Spatiotemporal} (\texttt{use\_spatiotemporal\_profile = T}):
    the full amplitude $A(\mathrm{transverse}, t)$ is interpolated from
    file at every time step. In 2D the data is $A(y, t)$; in 3D it is
    $A(\mathrm{tr1}, \mathrm{tr2}, t)$.
  \item \textbf{Static spatial} (\texttt{use\_spatiotemporal\_profile = F}):
    a time-independent spatial profile is interpolated onto the boundary
    once at start-up; the temporal envelope is set analytically in the
    deck via \texttt{t\_profile}.
\end{itemize}

In both modes the phase can also be read from file with
\texttt{use\_phase\_from\_file = T}, on the same grid as the amplitude.

\textbf{Note on defaults:} when \texttt{use\_custom\_profile = T} and
\texttt{use\_spatiotemporal\_profile} is not given, epoch2d defaults to the
spatiotemporal mode (\texttt{T}) while epoch3d defaults to the static
spatial mode (\texttt{F}). Setting the flag explicitly is recommended.

\section{EPOCH2D}

An example 2D \texttt{laser} block using all custom profile features:

\begin{lstlisting}[style=deck]
begin:laser
  boundary = x_min
  intensity_w_cm2 = 1e22
  lambda = 1.0 * micron
  pol = 0

  use_custom_profile = T
  use_spatiotemporal_profile = T
  profile_data_file = amplitude.dat

  n_t = 50
  n_y = 30
  y_min = -10.0 * micron
  y_max =  10.0 * micron
  t_start = 0.0
  t_end = 30e-15

  use_phase_from_file = T
  phase_data_file = phase.dat
end:laser
\end{lstlisting}

A 2D custom-profile laser block accepts the following parameters:

\begin{description}
  \item[\texttt{use\_custom\_profile}] Enable custom profile injection.
    Must be set to \texttt{T} to activate any of the parameters below. The
    default is \texttt{F}, which leaves all standard EPOCH profile
    behaviour unchanged.

  \item[\texttt{use\_spatiotemporal\_profile}] When \texttt{T}, the full
    two-dimensional $A(y, t)$ amplitude is read from a raw binary file and
    re-interpolated at every time step. When \texttt{F}, a one-dimensional
    static spatial profile $A(y)$ is read (also raw binary) and
    interpolated onto the boundary once at start-up; the temporal
    envelope is then set analytically via \texttt{t\_profile}. Defaults
    to \texttt{T} in epoch2d. 

  \item[\texttt{profile\_data\_file}] Path to the amplitude file. Relative
    paths resolve from the deck's data directory; absolute paths are used
    as-is. Defaults to \texttt{temporal\_spatial\_profile.dat}
    (spatiotemporal) or \texttt{spatial\_profile.dat} (static) if
    omitted.

  \item[\texttt{n\_t\_points} / \texttt{n\_t}] Number of time points in
    the profile array. Mandatory when
    \texttt{use\_spatiotemporal\_profile = T}.

  \item[\texttt{n\_transverse\_points} / \texttt{n\_y}] Number of
    transverse spatial points in the profile array ($\geq 2$). Mandatory
    whenever \texttt{use\_custom\_profile = T} (both modes).

  \item[\texttt{profile\_transverse\_min} / \texttt{y\_min}] Minimum
    transverse coordinate of the profile grid (metres). Mandatory
    whenever \texttt{use\_custom\_profile = T}. (For a laser on
    \texttt{y\_min}/\texttt{y\_max} the transverse axis is physically
    $x$; the \texttt{y\_min} alias name is kept for backward
    compatibility, but the boundary-agnostic
    \texttt{profile\_transverse\_min} reads better there.)

  \item[\texttt{profile\_transverse\_max} / \texttt{y\_max}] Maximum
    transverse coordinate of the profile grid (metres). Mandatory
    whenever \texttt{use\_custom\_profile = T}.

  \item[\texttt{t\_start}, \texttt{t\_end}] Start and end times of the
    profile grid (seconds). These reuse EPOCH's standard laser elements;
    the profile time axis runs from \texttt{t\_start} to \texttt{t\_end}
    with \texttt{n\_t} uniformly spaced points. Outside this window the
    laser is inactive.

  \item[\texttt{use\_phase\_from\_file}] When \texttt{T}, the phase is
    also read from file: on the spatiotemporal path $\varphi(y, t)$ is
    re-interpolated each time step; on the static path $\varphi(y)$ is
    interpolated once at start-up. Any \texttt{phase = ...} expression in
    the deck is ignored. Defaults to \texttt{F}, in which case the phase
    is set as normal via the \texttt{phase} deck element.

  \item[\texttt{phase\_data\_file}] Path to the raw binary phase file
    (radians). Follows the same path resolution as
    \texttt{profile\_data\_file}. The phase file must lie on the same
    grid as the amplitude file (same \texttt{n\_t}, \texttt{n\_y},
    bounds, \texttt{t\_start}, \texttt{t\_end}). Defaults to
    \texttt{phase\_profile.dat} if omitted.
\end{description}

\section{EPOCH3D}

In 3D the boundary plane has two transverse axes, so the grid declaration
uses two transverse point counts and two pairs of bounds. The
boundary-agnostic names are \texttt{tr1} and \texttt{tr2}, whose physical
meaning depends on the boundary the laser is attached to:

\begin{table}[h]
\centering
\begin{tabular}{@{}lll@{}}
\toprule
Boundary & tr1 axis & tr2 axis \\
\midrule
\texttt{x\_min}/\texttt{x\_max} & y & z \\
\texttt{y\_min}/\texttt{y\_max} & x & z \\
\texttt{z\_min}/\texttt{z\_max} & x & y \\
\bottomrule
\end{tabular}
\end{table}

An example 3D \texttt{laser} block (spatiotemporal, amplitude + phase, on
\texttt{x\_min}, written with the axis-named aliases):

\begin{lstlisting}[style=deck]
begin:laser
  boundary = x_min
  intensity_w_cm2 = 1e22
  lambda = 1.0 * micron
  pol = 0

  use_custom_profile = T
  use_spatiotemporal_profile = T
  profile_data_file = amplitude.dat

  n_y_points = 48
  n_z_points = 40
  n_t_points = 16
  y_min = -8.0 * micron
  y_max =  8.0 * micron
  z_min = -8.0 * micron
  z_max =  8.0 * micron
  t_start = 0.0
  t_end = 60e-15

  use_phase_from_file = T
  phase_data_file = phase.dat
end:laser
\end{lstlisting}

The same block written with the boundary-agnostic names (useful for decks
generated by boundary-independent tooling):

\begin{lstlisting}[style=deck]
  n_tr1 = 48
  n_tr2 = 40
  n_t = 16
  tr1_min = -8.0 * micron
  tr1_max =  8.0 * micron
  tr2_min = -8.0 * micron
  tr2_max =  8.0 * micron
\end{lstlisting}

A 3D custom-profile laser block accepts the following parameters (in
addition to \texttt{use\_custom\_profile},
\texttt{use\_spatiotemporal\_profile}, \texttt{profile\_data\_file},
\texttt{use\_phase\_from\_file}, \texttt{phase\_data\_file},
\texttt{t\_start} and \texttt{t\_end}, which behave exactly as in 2D):

\begin{description}
  \item[\texttt{n\_transverse1\_points} / \texttt{n\_tr1}] Number of
    points along the first transverse axis (see the table above).
    Mandatory ($\geq 2$) whenever \texttt{use\_custom\_profile = T}.

  \item[\texttt{n\_transverse2\_points} / \texttt{n\_tr2}] Number of
    points along the second transverse axis. Mandatory ($\geq 2$)
    whenever \texttt{use\_custom\_profile = T}.

  \item[\texttt{n\_t\_points} / \texttt{n\_t}] Number of time points.
    Mandatory ($\geq 2$) when \texttt{use\_spatiotemporal\_profile = T};
    ignored on the static spatial path.

  \item[\texttt{profile\_transverse1\_min} / \texttt{tr1\_min},
    \texttt{profile\_transverse1\_max} / \texttt{tr1\_max}] Bounds of the
    profile grid along the first transverse axis (metres). Mandatory
    whenever \texttt{use\_custom\_profile = T}.

  \item[\texttt{profile\_transverse2\_min} / \texttt{tr2\_min},
    \texttt{profile\_transverse2\_max} / \texttt{tr2\_max}] Bounds along
    the second transverse axis (metres). Mandatory whenever
    \texttt{use\_custom\_profile = T}.

  \item[\texttt{n\_\{x,y,z\}\_points} / \texttt{n\_\{x,y,z\}},
    \texttt{\{x,y,z\}\_min}, \texttt{\{x,y,z\}\_max}] Axis-named aliases
    for the point counts and bounds above. These are resolved to
    \texttt{tr1} or \texttt{tr2} according to the boundary table, so a
    deck for an \texttt{x\_min} laser can say
    \texttt{n\_y\_points}/\texttt{n\_z\_points} and
    \texttt{y\_min}...\texttt{z\_max} naturally. Naming the propagation
    axis (e.g.\ \texttt{n\_x\_points} on an \texttt{x\_min} laser) is an
    error --- EPOCH reports the offending element and line and suggests
    the transverse or boundary-agnostic names.
\end{description}

As in epoch2d, the static spatial path uses the raw binary format: the phase plane is
interpolated onto the boundary once at start-up alongside the amplitude.

\section{File Format}

All data files are raw binary with no embedded header. Shape and grid
bounds are declared entirely in the deck (see above). This matches
EPOCH's documented binary-file convention used elsewhere in the code
(e.g.\ \texttt{particles\_from\_file}). EPOCH checks the file size
against the declared shape at load time and aborts with the expected and
actual byte counts if they disagree. Note that the size check cannot
detect a transposed array of the same total size --- the axis ordering
below must be respected.

Each file contains an array of 64-bit IEEE~754 floating point values (the
same precision as EPOCH's \texttt{REAL(num)} type), written in
Fortran-compatible column-major order:

\begin{itemize}
  \item \textbf{2D spatiotemporal} $A(y, t)$: transverse $y$
    fastest-varying, time slowest. $n_y \times n_t$ values.
  \item \textbf{2D static spatial} $A(y)$: a single line of $n_y$ values.
  \item \textbf{3D spatiotemporal} $A(\mathrm{tr1}, \mathrm{tr2}, t)$:
    tr1 fastest-varying, then tr2, time slowest.
    $n_{\mathrm{tr1}} \times n_{\mathrm{tr2}} \times n_t$ values.
  \item \textbf{3D static spatial} $A(\mathrm{tr1}, \mathrm{tr2})$: tr1
    fastest-varying. $n_{\mathrm{tr1}} \times n_{\mathrm{tr2}}$ values.
\end{itemize}

All grids must be uniform; EPOCH reconstructs the axes from the declared
bounds and point counts (\texttt{np.linspace} on the Python side).
Amplitude values should be normalised to the range $[0, 1]$; EPOCH
multiplies by the peak field derived from \texttt{intensity\_w\_cm2} (or
\texttt{amp}). Phase values are in radians and are used directly ---
EPOCH adds no constant offset.

Interpolation is linear in every axis (bilinear for 2D spatiotemporal and
3D static, trilinear for 3D spatiotemporal). On the spatiotemporal path,
positions outside the declared transverse bounds return \textbf{zero}
amplitude and phase; on the static paths, positions outside the file grid
are \textbf{clamped} to the edge values. In either case the profile grid
should comfortably cover the boundary region the laser is meant to
illuminate.

\subsection{Laser profile generators}

The Fortran column-major layouts above correspond to NumPy arrays in
default C (row-major) order with the axes reversed, written with
\texttt{.tofile()} --- \textbf{no transpose is needed} in either
dimension:

\begin{table}[h]
\centering
\begin{tabular}{@{}ll@{}}
\toprule
File & NumPy shape (C order) \\
\midrule
2D spatiotemporal & $(n_t, n_y)$ \\
2D static spatial & $(n_y,)$ \\
3D spatiotemporal & $(n_t, n_{\mathrm{tr2}}, n_{\mathrm{tr1}})$ \\
3D static spatial & $(n_{\mathrm{tr2}}, n_{\mathrm{tr1}})$ \\
\bottomrule
\end{tabular}
\end{table}

2D example:

\begin{lstlisting}[style=python]
import numpy as np

N_T   = 50
N_Y   = 30
Y_MIN, Y_MAX   = -10e-6, 10e-6
T_START, T_END = 0.0, 30e-15

y = np.linspace(Y_MIN, Y_MAX, N_Y)
t = np.linspace(T_START, T_END, N_T)

# meshgrid with indexing="ij" gives arr[t_index, y_index]
Tg, Yg = np.meshgrid(t, y, indexing="ij")

# --- fill amplitude (values in [0, 1]) ---
t0, tau = 15e-15, 5e-15
w0 = 3e-6
amplitude = np.exp(-((Tg - t0) / tau)**2) * np.exp(-(Yg / w0)**2)

assert amplitude.shape == (N_T, N_Y)
amplitude.astype(np.float64).tofile("amplitude.dat")

# --- fill phase (values in radians) ---
phase = np.zeros((N_T, N_Y))  # flat phase for a plane wave
phase.astype(np.float64).tofile("phase.dat")
\end{lstlisting}

3D example (for an \texttt{x\_min} laser, so tr1 = y and tr2 = z):

\begin{lstlisting}[style=python]
import numpy as np

N_Y, N_Z, N_T  = 48, 40, 16          # n_tr1, n_tr2, n_t in the deck
Y_MIN, Y_MAX   = -8e-6, 8e-6
Z_MIN, Z_MAX   = -8e-6, 8e-6
T_START, T_END = 0.0, 60e-15

y = np.linspace(Y_MIN, Y_MAX, N_Y)
z = np.linspace(Z_MIN, Z_MAX, N_Z)
t = np.linspace(T_START, T_END, N_T)

# indexing="ij" gives arr[t_index, z_index, y_index] = (n_t, n_tr2, n_tr1)
Tg, Zg, Yg = np.meshgrid(t, z, y, indexing="ij")

t0, tau = 30e-15, 10e-15
w0 = 3e-6
amplitude = (np.exp(-((Tg - t0) / tau)**2)
             * np.exp(-(Yg**2 + Zg**2) / w0**2))

assert amplitude.shape == (N_T, N_Z, N_Y)
amplitude.astype(np.float64).tofile("amplitude.dat")

phase = np.zeros((N_T, N_Z, N_Y))
phase.astype(np.float64).tofile("phase.dat")
\end{lstlisting}

\section{Phase Convention}

EPOCH injects the field as:

\begin{equation}
E \propto \mathrm{profile} \cdot \sin(\omega t + \mathrm{phase})
\end{equation}

where $\omega t$ (the carrier phase) is accumulated internally by EPOCH.
The phase file must therefore contain only the slowly varying envelope
phase, not the carrier. This applies identically in 2D and 3D.

When using a complex envelope from LASY (where the physical field is
$E = \mathrm{Re}[E \cdot \exp(-i\omega_0 t)] = |E| \cdot \cos(\omega_0 t -
\varphi)$), the required conversion is:

\begin{equation}
\mathrm{phase\_epoch} = -(\mathrm{phase\_lasy} - \phi_{\mathrm{ref}}) -
\mathrm{gouy\_phase}
\end{equation}

where $\phi_{\mathrm{ref}} = \angle E(\mathrm{on\ axis}, t = t_{\mathrm{peak}})$
is the on-axis phase at peak time (which removes the arbitrary global
piston phase introduced by the propagator), and $\mathrm{gouy\_phase}$
pins the result to EPOCH's carrier convention. A $\pi/2$ offset converts
$\cos \to \sin$ but as a constant it does not affect field shape --- the
sign and reference subtraction are what matter physically.

\section{Two-Channel Polarisation}

A single custom-file laser block drives one transverse field component.
To inject a field with independent control of \emph{both} transverse
components at a boundary --- for example to synthesise an elliptically
or circularly polarised drive, or any polarisation state not reducible
to a single scalar profile --- add a second custom-file laser block on
the \emph{same} boundary with \texttt{pol} rotated by $90^{\circ}$
relative to the first. EPOCH's \texttt{pol} angle selects which pair of
transverse field components (e.g.\ \texttt{Ey} and \texttt{Ez} for an
\texttt{x\_min} laser) a laser block drives, so \texttt{pol = 0} and
\texttt{pol = 90} address the two components independently, each from
its own file:

\begin{lstlisting}[style=deck]
begin:laser
  boundary = x_min
  ...
  pol = 0
  profile_data_file = amplitude_H.dat   # drives Ey
  ...
end:laser

begin:laser
  boundary = x_min
  ...
  pol = 90
  profile_data_file = amplitude_V.dat   # drives Ez
  ...
end:laser
\end{lstlisting}

The loaded amplitude and phase data live on the \texttt{laser\_block}
instance itself rather than in a shared module-level variable, so each
block keeps its own copy: the second block's file cannot silently
overwrite or alias the first's, even when the two grids differ in shape
or resolution. (Prior to this per-laser storage, a second custom-file
block on the same boundary would incorrectly reuse the first block's
already-loaded data.) The \texttt{cos}/\texttt{sin(pol\_angle)} channel
split that routes each block to its transverse component is identical in
every boundary routine of epoch2d and epoch3d, so the same two-block
pattern works unchanged in 3D. (The two-channel configuration has been
validated end-to-end in 2D; the 3D per-laser storage is the same design,
built in from the start.)

\section{Current Limitations}

\begin{itemize}
  \item \textbf{epoch1d} has no custom-profile support.
  \item All file grids must be \textbf{uniform}; non-uniform coordinate
    axes are not supported.
  \item Memory: in \textbf{3D}, each MPI rank \emph{can} store only the
    slab of the spatiotemporal array covering its own patch of the laser
    boundary (plus a one-cell margin) rather than the full plane, with
    loading streamed one time-slice at a time. This is now the default
    behaviour for a static domain decomposition: \texttt{use\_pre\_balance}
    defaults to \texttt{T} in epoch3d and triggers a one-off startup load
    balance that can move a rank's local patch, so the file load for a
    spatiotemporal laser is deferred past that rebalance and windowed
    against the settled domain instead of being sized too early.
    Per-rank cost is then roughly $(\text{local patch fraction}) \times
    n_{\mathrm{tr1}} \times n_{\mathrm{tr2}} \times n_t \times 8$ bytes.
    Continuous dynamic load balancing (\texttt{use\_balance = T}) keeps
    the same per-rank saving for the whole run rather than only until
    the first rebalance: the window is re-derived and the file reloaded
    after every genuine mid-run redistribution, at the same per-event
    cost as the initial load.

    The full plane is still stored on every rank, unconditionally, for
    a \emph{moving injection window} on a \texttt{y}/\texttt{z}-boundary
    laser (its tr1 axis is $x$, which the window shifts as the
    simulation window moves). This is deliberately retained rather than
    windowed: unlike domain rebalancing, the window's motion is not
    tied to a discrete, all-ranks-agree event such an update could hook
    into, so a safe local-slab alternative does not exist yet for this
    case. In \textbf{2D} the (small) full array is always broadcast to
    every rank.
\end{itemize}

\end{document}
